\section{The theta kernel and the original closed energy}\label{sec:energy}
In the foundation proofs write $s_n=\ell_n=\log n$.

For $x\ge0$, define
\begin{equation}\label{eq:kappa}
 \kappa(x)=e^{x/2}\sum_{j\ge1}(4\tau_j(x)^2-6\tau_j(x))e^{-\tau_j(x)},
 \qquad \tau_j(x)=\pi j^2e^{2x},
\end{equation}
and extend evenly. We record the normalization and decay needed later.

\begin{lemma}[Smoothness and normalization]\label{lem:theta}
The kernel $\kappa$ is smooth, strictly positive, and even. There is a finite $C_\kappa>0$ such that
\begin{equation}\label{eq:envelopek}
 \kappa(x)\le C_\kappa\exp\!\left(\frac92|x|-\pi e^{2|x|}\right),
 \qquad
 C_\kappa=4\pi^2\sum_{j\ge1}j^4e^{-\pi(j^2-1)}.
\end{equation}
Furthermore $\int_\R2c\kappa\,dx=1$, and
\begin{equation}\label{eq:envelopeW}
 0<W(x)\le C_\kappa\exp(4|x|-\pi e^{2|x|}).
\end{equation}
\end{lemma}
\begin{proof}
The series in~\eqref{eq:kappa}, together with every derivative, converges uniformly on compact subsets of $[0,\infty)$. Each summand is positive because $\tau_j\ge\pi>3/2$. Dropping the negative polynomial term and using $e^{2x}\ge1$ gives~\eqref{eq:envelopek} for $x\ge0$.

To justify smooth evenness at zero, rather than assume it from reflection, set
\[
 H(x)=e^{x/2}\sum_{j\ge1}e^{-\pi j^2e^{2x}},\qquad x\in\R.
\]
The Jacobi transformation \cite[Eq.~20.7.32]{DLMF} implies
$H(-x)=H(x)+\sinh(x/2)$. Applying $D^2-1/4$ shows that $H''-H/4$ is even; direct differentiation identifies it with~\eqref{eq:kappa} on the positive half-line. Thus it gives the claimed smooth extension. Differentiating the transformation at zero yields $H'(0)=-1/4$. Since $c''=c/4$, integration by parts and the rapid decay of $H,H'$ at $+\infty$ give
\[
 \int_0^\infty c(x)\kappa(x)\,dx
  =[cH'-c'H]_0^\infty=\frac14.
\]
Evenness proves the normalization. Finally $2c(x)\ge e^{|x|/2}$ gives~\eqref{eq:envelopeW}.
\end{proof}

Put $\rho(s)=e^{-s/2}/(1-e^{-2s})$ for $s>0$, and $a_n=\Lambda(n)/\sqrt n$. The full positive energy is
\begin{align}
 \E_\Gamma[f]&=\int_0^\infty\rho(s)\int_\R
       \kappa(x)\kappa(x+s)|f(x+s)-f(x)|^2\,dx\,ds,\label{eq:Eg}\\
 \E_P[f]&=\sum_{n\ge2}a_n\int_\R
       \kappa(x)\kappa(x+s_n)|f(x+s_n)-f(x)|^2\,dx,\label{eq:Ep}\\
 \E[f]&=\E_\Gamma[f]+\E_P[f],\qquad
 \D=\{f\in\He:\E[f]<\infty\}.\label{eq:Edomain}
\end{align}
These are nonnegative integrals and sums, interpreted by Tonelli's theorem; no cancellation defines their domain. Write
\[
 \mu(f)=\int_\R f\,d\mu,\qquad
 \Var_\mu(f)=\norm{f}_\mu^2-|\mu(f)|^2.
\]

\begin{proposition}[Maximal domain and compact smooth form core]\label{prop:core}
The form $\E$ on $\He$ is densely defined, nonnegative, and closed. Its maximal integral domain $\D$ is the closure of even $C_c^\infty(\R)$ in the norm $(\norm{f}_\mu^2+\E[f])^{1/2}$. Constants belong to $\D$, and $\E[f]=0$ holds precisely for almost-everywhere constant $f$.
\end{proposition}
\begin{proof}
The theta envelope implies that for every $A_0>0$ there is a constant $C_{A_0}$ such that $\kappa(x)\le C_{A_0}e^{-A_0|x|}$. Consequently
\begin{equation}\label{eq:edgefinite}
 \int_\R\kappa(x)\kappa(x+s)\,dx
 \le C_{A_0}^2(s+A_0^{-1})e^{-A_0s},\qquad s\ge0.
\end{equation}
Since $a_n\le n$, the complete prime edge measure has finite mass, by taking $A_0>2$ in~\eqref{eq:edgefinite}. For every $L<\infty$,
\begin{equation}\label{eq:rhointegrable}
 \int_0^\infty\rho(s)\min\{1,L^2s^2\}\,ds<\infty.
\end{equation}
These observations admit compact smooth functions and constants to the domain.

For closedness, let $(f_j)$ be Cauchy in the form norm, with $L^2(\mu)$ limit $f$. A subsequence converges almost everywhere for $\mu$, hence for Lebesgue measure, because the density is strictly positive. A countable union of translated null sets is null, so all prime differences converge almost everywhere. Fubini's theorem gives convergence for the Gamma edge measure as well. For each fixed $k$, Fatou's lemma gives
\[
 \E[f-f_k]\le\liminf_{j\to\infty}\E[f_j-f_k].
\]
It follows that $f\in\D$ and $f_k\to f$ in form norm.

We prove the core assertion first for real profiles. Let $T_M(t)=\max\{-M,\min\{t,M\}\}$. The remainder $r_M(t)=t-T_M(t)$ is $1$-Lipschitz, tends pointwise to zero, and satisfies $|r_M(t)|\le|t|$. Thus dominated convergence in both edge measures and in $L^2(\mu)$ gives $T_M(f)\to f$ in form norm. It suffices to approximate bounded even profiles.

Choose smooth even cutoffs $0\le\chi_R\le1$ equal to one on $[-R,R]$, supported in $[-2R,2R]$, and with uniformly bounded Lipschitz constants. Equation~\eqref{eq:rhointegrable}, the finite prime-edge mass, and dominated convergence give $\E[\chi_R]\to0$. For bounded $f$, the identity for a product difference yields
\begin{align*}
 |[f(1-\chi_R)](y)-[f(1-\chi_R)](x)|^2
 &\le2(1-\chi_R(y))^2|f(y)-f(x)|^2\\
 &\quad+2\norm{f}_\infty^2|\chi_R(y)-\chi_R(x)|^2.
\end{align*}
Hence $f\chi_R\to f$ in form norm. It remains to smooth a compactly supported profile $g\in\D$.

Suppose $\supp g\subset[-K,K]$. Its ordinary $L^2$ norm is finite, since the density of $\mu$ is bounded below on this interval. For $0<s\le1$, a nonzero difference of $g$ has both endpoints in $[-K-1,K+1]$. The positive lower bound of $\kappa$ there controls the small-shift part of $\tg[g]$ by $\E_\Gamma[g]$. The part with $s>1$ is bounded by $4\norm{g}_2^2\int_1^\infty\rho(s)\,ds$. Thus $\tg[g]<\infty$.

Let $g_\epsilon=\zeta_\epsilon*g$, where $\zeta_\epsilon$ is an even compactly supported smooth approximate identity. Translations commute with convolution, and
\[
 \norm{\Delta_s(g_\epsilon-g)}_2\le2\norm{\Delta_sg}_2,
 \qquad\Delta_sg=g(\cdot+s)-g.
\]
Dominated convergence gives $\tg[g_\epsilon-g]\to0$. Since $\kappa$ is bounded,
$\E_\Gamma[g_\epsilon-g]\le\norm{\kappa}_\infty^2\tg[g_\epsilon-g]\to0$.
All the smoothed profiles have support in one fixed compact interval $J$. For any $h$ supported there,
\begin{equation}\label{eq:compactprimerate}
 \E_P[h]\le2\int_J\kappa(x)|h(x)|^2
     \sum_{n\ge2}a_n\bigl(\kappa(x-s_n)+\kappa(x+s_n)\bigr)\,dx
 \le C_J\norm{h}_2^2.
\end{equation}
Indeed, the complete rate in parentheses is uniformly bounded on $J$: use $\kappa(x\pm\log n)\le C_{A_0}e^{A_0\sup_J|x|}n^{-A_0}$ and $A_0>2$. Ordinary $L^2$ convergence therefore proves prime-energy convergence. The local upper bound of the density gives $L^2(\mu)$ convergence. Applying the argument to real and imaginary parts proves the complex core assertion.

Finally $\E[1]=0$. If $\E[f]=0$, the strictly positive Gamma density implies $f(y)=f(x)$ for almost every pair $x<y$. Fubini's theorem forces $f$ to be almost everywhere constant. Density of the core in $\He$ is immediate from the positive smooth density of $\mu$.
\end{proof}

\section{The explicit formula and the compact physical bridge}\label{sec:bridge}

We now prove the identity that connects the positive energy to the physical form. The classical explicit formula is used as a theorem, without any assumption about the location of nontrivial zeta zeros inside the critical strip. We first fix its normalization, then justify its use with the noncompact theta kernel, and finally perform the ground-state calculation with all endpoints retained.

In this section use the nonunitary Fourier transform
\[
 \mathcal Fu(z)=\int_\R u(x)e^{izx}\,dx,
 \qquad m(u)=\int_\R c(x)u(x)\,dx.
\]
Let $\rho_\zeta$ range over the distinct nontrivial zeros of $\zeta$, with multiplicity $\nu_{\rho_\zeta}$, and put $w_{\rho_\zeta}=(\rho_\zeta-1/2)/i$. Thus $|\operatorname{Im}w_{\rho_\zeta}|<1/2$; no reality of $w_{\rho_\zeta}$ is asserted. We use the classical zero-counting bound
\begin{equation}\label{eq:zerocount}
 \sum_{|\operatorname{Im}\rho_\zeta|\le T}\nu_{\rho_\zeta}
       =O(T\log(T+2))\qquad(T\ge2).
\end{equation}
This weak form of the classical Riemann--von Mangoldt estimate \cite[Theorem~0.1]{GarrettZeros} is sufficient for every zero sum below.

\begin{theorem}[Classical Guinand explicit formula]\label{thm:guinand}
For $g\in C_c^\infty(\R)$,
\begin{align}\label{eq:guinand}
 \sum_{\rho_\zeta}\nu_{\rho_\zeta}\mathcal Fg(w_{\rho_\zeta})
 &=\mathcal Fg(i/2)+\mathcal Fg(-i/2)\notag\\
 &\quad+\frac1{2\pi}\int_\R
   \bigl(\operatorname{Re}\psi(1/4+it/2)-\log\pi\bigr)
       \mathcal Fg(t)\,dt\notag\\
 &\quad-\sum_{n\ge2}a_n\bigl(g(s_n)+g(-s_n)\bigr).
\end{align}
The zero sum is absolutely convergent. This is the classical formula in the Fourier convention of Garrett \cite[Theorem~0.1]{Garrett}, with prime powers collected by the von Mangoldt function.
\end{theorem}

The statement is quoted, not proposed as a new result. Absolute convergence follows from the rapid decay of $\mathcal Fg$ in the closed strip $|\operatorname{Im}z|\le1/2$ and~\eqref{eq:zerocount}.

\begin{lemma}[Bilinear physical normalization]\label{lem:bilinear}
For real even $u,v\in C_c^\infty(\R)$,
\begin{equation}\label{eq:bilinearzeros}
 q_0(u,v)+2m(u)m(v)
   =\sum_{\rho_\zeta}\nu_{\rho_\zeta}
       \mathcal Fu(w_{\rho_\zeta})\mathcal Fv(w_{\rho_\zeta}).
\end{equation}
\end{lemma}
\begin{proof}
The digamma integral \cite[Eq.~5.9.16]{DLMF} gives, after subtracting its value at $1/4$ and setting its integration variable equal to $2s$,
\begin{equation}\label{eq:digammamultiplier}
 \operatorname{Re}\psi(1/4+it/2)-\psi(1/4)
    =2\int_0^\infty\rho(s)(1-\cos(ts))\,ds=:d_\Gamma(t).
\end{equation}
The difference integral is convergent at zero and infinity. Plancherel and polarization identify
\[
 \tg(u,v)=\frac1{2\pi}\int_\R
       d_\Gamma(t)\mathcal Fu(t)\mathcal Fv(t)\,dt
\]
for real even tests. Apply Theorem~\ref{thm:guinand} to $g=u*v$. Its Fourier transform is the product of the two transforms, and $g(0)=\int uv$. Evenness gives
\[
 \mathcal Fu(i/2)=\mathcal Fu(-i/2)=m(u),\qquad
 g(s_n)+g(-s_n)
  =\int_\R\bigl(u(x+s_n)v(x)+u(x)v(x+s_n)\bigr)\,dx.
\]
These are respectively the pole and the polarized prime term. The constant in the Gamma term is $\psi(1/4)-\log\pi=-\cg$, proving~\eqref{eq:bilinearzeros}.
\end{proof}

For complex even tests the zero-side product is
$\mathcal Fu(w)\overline{\mathcal Fv(\overline w)}$. It is not generally $|\mathcal Fu(w)|^2$ when $u=v$ and $w$ is nonreal. We need only the real bilinear identity and its conjugate-linear extension in the second argument; no positivity is inferred from this zero sum.

\begin{lemma}[Entire theta transform]\label{lem:thetaxi}
Let
\[
 \xi(s)=\tfrac12s(s-1)\pi^{-s/2}\Gamma(s/2)\zeta(s).
\]
Then the entire continuation of this function satisfies
\begin{equation}\label{eq:thetaxi}
 \mathcal F\kappa(z)=\xi(1/2+iz),\qquad z\in\C.
\end{equation}
In particular, $\mathcal F\kappa(w_{\rho_\zeta})=0$ for every nontrivial zero, whether or not it is on the critical line.
\end{lemma}
\begin{proof}
Write $\Theta(t)=\sum_{j\ge1}e^{-\pi j^2t}$, so that the function in Lemma~\ref{lem:theta} is $H(x)=e^{x/2}\Theta(e^{2x})$. For $\operatorname{Re}s>1$, termwise integration gives
\[
 \pi^{-s/2}\Gamma(s/2)\zeta(s)
   =\int_0^\infty\Theta(t)t^{s/2-1}\,dt.
\]
The Jacobi transformation is
$\Theta(t)=t^{-1/2}\Theta(1/t)+(t^{-1/2}-1)/2$.
Split the integral at one and substitute $t\mapsto1/t$ in its lower part. The elementary remainder integrates to $1/[s(s-1)]$, and hence
\begin{align}\label{eq:mellintheta}
 \pi^{-s/2}\Gamma(s/2)\zeta(s)
 &=\frac1{s(s-1)}+I(s),\notag\\
 I(s)&=\int_1^\infty\Theta(t)
       \bigl(t^{s/2-1}+t^{(1-s)/2-1}\bigr)\,dt.
\end{align}
The integral $I$ is entire by local uniform double-exponential domination after $t=e^{2x}$. Thus
$\xi(s)=1/2+s(s-1)I(s)/2$ is entire. For $s=1/2+iz$, the same substitution gives
\begin{equation}\label{eq:xiH}
 \xi(1/2+iz)=\frac12
       -2(z^2+1/4)\int_0^\infty H(x)\cos(zx)\,dx.
\end{equation}
Because $\kappa=H''-H/4$ on the positive half-line, evenness and two integrations by parts give
\[
 \mathcal F\kappa(z)
  =-2H'(0)-2(z^2+1/4)\int_0^\infty H(x)\cos(zx)\,dx.
\]
All boundary terms at infinity vanish for each complex $z$. The value $H'(0)=-1/4$ from Lemma~\ref{lem:theta} proves~\eqref{eq:thetaxi}. The defining prefactors of $\xi$ are finite and nonzero at a nontrivial zero of $\zeta$, so the last assertion follows.
\end{proof}

\begin{proposition}[Weak theta source identity]\label{prop:thetanull}
For every even $h\in C_c^\infty(\R)$,
\begin{equation}\label{eq:thetanull}
 q_0(\kappa,h)=-\overline{m(h)}.
\end{equation}
All terms on the left exist separately as convergent form pairings. Equivalently, $(D_\Gamma-\cg-P)\kappa=-c$ on even compact smooth tests. This is a weak identity and does not assert that $c$ belongs to ordinary $L^2$.
\end{proposition}
\begin{proof}
We first take $h$ real. Let $\chi_R(x)=\chi(x/R)$, $R\ge1$, where $\chi$ is even and smooth, equals one on $[-1,1]$, has support in $[-2,2]$, and satisfies $0\le\chi\le1$. Put $k_R=\kappa\chi_R$. Differentiating the theta series on each half-line shows that $\kappa$ and its derivatives decrease faster than every exponential. Consequently $k_R\to\kappa$ in $H^1(\R)$ and in $L^1(e^{|x|/2}dx)$. For $v\in H^1(\R)$,
\[
 \tg[v]\le\norm{v'}_2^2\int_0^1\rho(s)s^2\,ds
              +4\norm{v}_2^2\int_1^\infty\rho(s)\,ds.
\]
It follows that the Gamma pairing, the ordinary $L^2$ pairing, and the moment in~\eqref{eq:bilinearzeros} converge as $R\to\infty$.

The complete prime pairing also converges without a spatially uniform incoming-rate assumption. If $\supp h\subset[-K,K]$, changing variables in one of its two terms gives
\[
 p(k_R,h)=\sum_{n\ge2}a_n\int_{-K}^K h(x)
           \bigl(k_R(x+s_n)+k_R(x-s_n)\bigr)\,dx.
\]
For any fixed $A_0>2$,
\begin{equation}\label{eq:bridgeprimerate}
 \sup_{|x|\le K}\sum_{n\ge2}a_n
       \bigl(\kappa(x+s_n)+\kappa(x-s_n)\bigr)
 \le2C_{A_0}e^{A_0K}\sum_{n\ge2}a_nn^{-A_0}<\infty.
\end{equation}
Here $a_n\le n$ suffices. Since $|k_R|\le\kappa$, dominated convergence proves the assertion and absolute convergence of the limiting pairing.

For the zero side, uniformly in $R$ and in $|\operatorname{Im}z|\le1/2$,
\[
 |\mathcal Fk_R(z)|\le\int_\R\kappa(x)e^{|x|/2}\,dx<\infty.
\]
Repeated integration by parts in the transform of the fixed compact test gives
$|\mathcal Fh(z)|\le C_{h,N}(1+|\operatorname{Re}z|)^{-N}$ throughout this strip. Equation~\eqref{eq:zerocount} therefore supplies a summable majorant for the products in~\eqref{eq:bilinearzeros}, for example with $N=3$. Pointwise transform convergence and Lemma~\ref{lem:thetaxi} make their limiting sum zero. Thus
\[
 q_0(\kappa,h)+2m(\kappa)m(h)=0.
\]
Lemma~\ref{lem:theta} gives $m(\kappa)=1/2$, proving~\eqref{eq:thetanull} for real $h$. Splitting an arbitrary $h$ into real and imaginary parts and using conjugate linearity in the second argument proves the stated version.
\end{proof}

\begin{theorem}[Compact physical bridge]\label{thm:compactbridge}
For every even $f\in C_c^\infty(\R)$,
\begin{equation}\label{eq:provedbridge}
 q_0[\kappa f]=\E[f]-\frac12\norm{f}_\mu^2.
\end{equation}
\end{theorem}
\begin{proof}
Set $h=\kappa|f|^2$, which is a real even compact smooth test. For $a=\kappa(x)$, $b=\kappa(x+s)$, $F=f(x)$, and $G=f(x+s)$, the pointwise identity
\[
 |bG-aF|^2-(b-a)(b|G|^2-a|F|^2)=ab|G-F|^2
\]
gives, after integration,
\begin{equation}\label{eq:gammaground}
 \tg[\kappa f]-\tg(\kappa,h)=\E_\Gamma[f].
\end{equation}
This is an equality between finite quantities: $\kappa$, $h$, and $\kappa f$ are in $H^1$, their Gamma pairings are finite by the estimate in the preceding proof, and $\E_\Gamma[f]$ is finite by Proposition~\ref{prop:core}.

For the primes the two endpoint terms are both essential. Their full expansion gives
\begin{align}\label{eq:primeground}
 p[\kappa f]-p(\kappa,h)
 &=\sum_{n\ge2}a_n\int_\R\kappa(x)\kappa(x+s_n)
       \bigl(2\operatorname{Re}(f(x+s_n)\overline{f(x)})\notag\\
 &\hspace{44mm}-|f(x+s_n)|^2-|f(x)|^2\bigr)\,dx\notag\\
 &=-\E_P[f].
\end{align}
The compact-profile prime energy is finite, the compact correlation $p[\kappa f]$ has only finitely many nonzero shifts, and $p(\kappa,h)$ is absolutely convergent by~\eqref{eq:bridgeprimerate}. Thus no difference of divergent prime expressions is used. Finally,
$\norm{\kappa f}_2^2=\ip{\kappa}{h}_2$ cancels the constant term. Combining~\eqref{eq:gammaground}, \eqref{eq:primeground}, and~\eqref{eq:thetanull} yields
\[
 q_0[\kappa f]=\E[f]-m(h)
             =\E[f]-\int_\R c\kappa|f|^2\,dx,
\]
which is~\eqref{eq:provedbridge}.
\end{proof}

The theorem uses the classical explicit formula and the theta transform, not a zero-location hypothesis. It establishes an identity of forms, not a positivity statement for $q_0$ or for the centered energy shifted by $1/2$.

\section{The complete weighted prime form}\label{sec:prime}
In this appendix $\norm u_\omega^2=\int_\R |u|^2/W\,dx$ is the full-line weighted numerator norm.

The relevant prime bound is in the numerator's weighted norm, not a uniform vanishing estimate for an incoming jump rate. Let $\beta=\sqrt W$ and let $\tau_sh(x)=h(x+s)$ on ordinary $L^2(\R)$.

\begin{lemma}[Weighted shift estimate]\label{lem:prime}
For a finite constant $K$ and every $s\ge0$,
\begin{equation}\label{eq:betaprod}
 \sup_{x\in\R}\beta(x)\beta(x+s)\le K e^{2s-\pi e^s}.
\end{equation}
The self-adjoint series
\begin{equation}\label{eq:KP}
 K_P=\sum_{n\ge2}a_nM_\beta(\tau_{s_n}+\tau_{-s_n})M_\beta
\end{equation}
converges absolutely in operator norm and preserves even functions. If $K_{P,N}$ denotes its head through $n=N$, then
\begin{equation}\label{eq:rN}
 \norm{K_P-K_{P,N}}\le r_N,
 \qquad r_N=2K\sum_{n>N}\Lambda(n)n^{3/2}e^{-\pi n}\longrightarrow0.
\end{equation}
For every $u$ of finite weighted norm the series~\eqref{eq:physicalprime} converges absolutely and
\begin{align}
 p[u]&=\ip{K_P(u/\beta)}{u/\beta}_2,\label{eq:prepresentation}\\
 |p[u]-p_N[u]|&\le r_N\norm{u}_\omega^2,\qquad
 |p_N[u]|\le2\sum_{2\le n\le N}a_n\norm{u}_2^2.\label{eq:headtail}
\end{align}
\end{lemma}
\begin{proof}
Set $z=|x|+|x+s|\ge s$. Equation~\eqref{eq:envelopeW} and the arithmetic--geometric mean inequality give
\begin{align*}
 \beta(x)\beta(x+s)
 &\le K\exp\!\left(2z-\frac\pi2(e^{2|x|}+e^{2|x+s|})\right)\\
 &\le K e^{2z-\pi e^z}\le K e^{2s-\pi e^s}.
\end{align*}
The last step uses $2-\pi e^z<0$ for $z\ge0$. The norm of each weighted translation is bounded by the supremum of its multiplier, so its two-shift contribution at $s_n$ has norm at most $2Ka_nn^2e^{-\pi n}$. Since $\Lambda(n)\le\log n$, the norm series converges. Adjointness interchanges the two shifts, and reflection also interchanges them. This proves self-adjointness and even invariance.

Writing $u=\beta F$, the $n$th quadratic term equals the corresponding term in~\eqref{eq:physicalprime}. Its absolute value is bounded by the summable operator-norm bound times $\norm{F}_2^2=\norm{u}_\omega^2$. This proves both the identification and absolute convergence. The tail estimate follows from~\eqref{eq:rN}; Cauchy--Schwarz and translation invariance of ordinary $L^2$ give the finite-head estimate.
\end{proof}

\begin{remark}
Lemma~\ref{lem:prime} does not assert that $K_P$ is compact. Multiplication by a decaying function followed by a translation need not be compact on a nonatomic $L^2$ space. Nor does~\eqref{eq:rN} assert a uniformly vanishing tail for the unweighted incoming prime rate as the physical point escapes to infinity. The norm, the weight, and the coordinates in the lemma are essential.
\end{remark}

\section{The exact numerator domain}\label{sec:domain}
Define
\begin{equation}\label{eq:X}
 \begin{gathered}
 X=\{u\in\Le:\tg[u]<\infty,\ \norm u_\omega<\infty\},\\
 \norm u_\omega^2=\int_\R |u|^2/W\,dx,\qquad
 \norm u_X^2=\tg[u]+\norm u_\omega^2.
 \end{gathered}
\end{equation}

We first analyze $X$ directly from its two defining nonnegative terms. Because $W$ is positive and bounded,
\begin{equation}\label{eq:L2bound}
 \norm{u}_2^2\le\norm{W}_\infty\norm{u}_\omega^2.
\end{equation}
The Gamma form is closed on ordinary $L^2$, for example by the same difference-and-Fatou argument used for $\E$.

\begin{lemma}\label{lem:Xcore}
The space $X$ of~\eqref{eq:X} is Hilbert, and even $C_c^\infty(\R)$ is dense in $X$.
\end{lemma}
\begin{proof}
A sequence Cauchy in $X$ has a weighted $L^2$ limit and an ordinary Gamma-form limit. Equation~\eqref{eq:L2bound} identifies their ordinary $L^2$ limits, proving completeness.

Use the cutoffs $\chi_R$ from Proposition~\ref{prop:core}. For $h_R=(1-\chi_R)u$ the product identity gives
\begin{align*}
 |h_R(x+s)-h_R(x)|^2
 &\le2(1-\chi_R(x+s))^2|u(x+s)-u(x)|^2\\
 &\quad+2|u(x)|^2|\chi_R(x+s)-\chi_R(x)|^2.
\end{align*}
The first term tends to zero in the Gamma integral by dominated convergence. The second is dominated by
$2|u(x)|^2\min\{1,L^2s^2\}$, which is integrable by~\eqref{eq:rhointegrable}. Weighted convergence is another application of dominated convergence. Thus compactly supported vectors are dense in $X$.

For a fixed compactly supported vector, even smooth mollification converges in ordinary $L^2$ and in $\tg$, by the translation argument in Proposition~\ref{prop:core}. The supports remain in one compact interval, where $1/W$ is bounded. Ordinary $L^2$ convergence consequently gives weighted convergence as well.
\end{proof}

For original profiles and physical numerators the coordinates are
\begin{equation}\label{eq:coordinates}
 u=\kappa f,\qquad v=2cf=\frac uW,\qquad
 F=\frac u{\sqrt W}=\sqrt{2c\kappa}\,f.
\end{equation}
Only $f\mapsto F$ is the displayed unitary map from $\He$ to ordinary $\Le$. The map $f\mapsto u$ is an isometry to the weighted space with norm $\norm{\cdot}_\omega$, since
\begin{equation}\label{eq:weightedidentity}
 \norm{\kappa f}_\omega=\norm{f}_\mu.
\end{equation}

\begin{proposition}[Extension of the compact bridge]\label{prop:domain}
Multiplication by $\kappa$ gives a bijection $\D\to X$, with equivalent graph norms, and
\begin{equation}\label{eq:extendedbridge}
 q_0(\kappa f,\kappa g)=\E(f,g)-\tfrac12\ip{f}{g}_\mu,
 \qquad f,g\in\D.
\end{equation}
\end{proposition}
\begin{proof}
Put $C_0=|\cg|\norm{W}_\infty+\norm{K_P}$. On the compact smooth core, Theorem~\ref{thm:compactbridge}, Lemma~\ref{lem:prime}, and~\eqref{eq:weightedidentity} imply
\begin{align}
 \tg[\kappa f]&\le\E[f]+(C_0+\tfrac12)\norm{f}_\mu^2,\label{eq:graph1}\\
 \E[f]&\le\tg[\kappa f]+(C_0+\tfrac12)\norm{f}_\mu^2.\label{eq:graph2}
\end{align}
Taking compact-profile approximants from Proposition~\ref{prop:core} and using completeness of $X$ proves that $\kappa\D\subset X$.

Conversely, for $u\in X$ choose even $u_j\in C_c^\infty$ converging in $X$, by Lemma~\ref{lem:Xcore}. Since $\kappa$ is smooth and strictly positive, $f_j=u_j/\kappa$ are legitimate even compact smooth profiles. Equation~\eqref{eq:graph2} makes $(f_j)$ Cauchy in the original energy graph norm. Its weighted limit is $f=u/\kappa$, so closedness of $\E$ gives $f\in\D$. This proves surjectivity and graph-norm equivalence. The Gamma form and the prime form are continuous in $X$, so the quadratic identity passes to the limit; polarization gives~\eqref{eq:extendedbridge}.
\end{proof}

Recall $m(u)=\int_\R c(x)u(x)\,dx$. The integral exists absolutely for finite weighted norm, because
\begin{equation}\label{eq:meanbound}
 |m(u)|\le\left(\int_\R c^2W\,dx\right)^{1/2}\norm{u}_\omega
       =\tfrac12\norm{u}_\omega.
\end{equation}
Here $\int c^2W=\tfrac12\int c\kappa=1/4$. Thus $\mu(f)=2m(\kappa f)$, and the full pole form obeys
\begin{equation}\label{eq:polebridge}
 q[u]:=q_0[u]+2|m(u)|^2
       =\E[f]-\tfrac12\Var_\mu(f),\qquad u=\kappa f\in X.
\end{equation}
The coefficient $2$ in the pole is fixed by normalization. It is not an additional positive contribution available after centering.


\section{Sharp cuts and the centered exterior minimum domain}\label{app:cuts}
We record the domain argument, rather than identifying a minimum
closure with an arbitrary maximal finite-energy class.
Write $W_\R(h)=\int_\R |h|^2/\alpha^2$. Under $h=\kappa f$ it equals
$\norm f_{L^2(\mu)}^2$. Proposition~\ref{prop:domain}
shows that the original topology and
$(t_\Gamma(h)+W_\R(h))^{1/2}$ are equivalent on numerators.

For a smooth numerator $h$, expansion gives
\[
 t_\Gamma(h1_I)+t_\Gamma(h1_O)
 =t_\Gamma(h)+2\ip{h1_I}{\J_\Gamma(h1_O)}.
\]
For $0<s\le2$, \eqref{eq:rhoelementary} gives
$\rho(s)\le1/(2s)+1/4\le1/s$. For $s\ge2$, the function
$se^{-s/2}$ decreases, so
$s\rho(s)\le(2/e)/(1-e^{-4})<1$; the last inequality follows already
from $e>8/3$. Thus the unweighted Carleman bound for the Gamma crossing
is finite: $\rho(s)\le1/s$ and the two boundary Carleman maps give,
for example, $\norm{\J_\Gamma}_{L^2(O)\to L^2(I)}\le2\pi$.
Thus the right-hand correction is bounded by a constant times
$\norm h_2^2$, and $W_\R$ splits exactly. Since $\alpha^2$ is bounded,
both cut maps extend continuously in the numerator form topology.

For a fixed smooth compact numerator, approximate either step cut from
its own side by smooth killed-core cutoffs. The error $e_\delta$ is
supported in strips of width $\delta$ at $\pm B$ and satisfies
\[
 \norm{e_\delta}_\infty+\operatorname{TV}(e_\delta)\le C,
 \qquad\norm{e_\delta}_2^2+W_\R(e_\delta)\le C\delta.
\]
The bounded-variation translation estimate gives
\[
 \norm{e_\delta(\cdot+s)-e_\delta}_2^2\le C\min(s,\delta).
\]
Integrating against $\rho(s)=O(1/s)$ near zero and its integrable
large-shift tail proves
$t_\Gamma(e_\delta)=O(\delta(1+|\log\delta|))\to0$.
Consequently the cuts are in the corresponding killed minimum closures.
Apply continuity to a general minimum-domain numerator to obtain the
same conclusion. On the bounded local interval $W_\R$ is equivalent
to physical $L^2$, because $\alpha$ is continuous and strictly positive.

The constant function belongs to $\D$. Indeed even smooth cutoffs
$\chi_R$ of the constant satisfy $\kappa\chi_R\to\kappa$ in physical
$H^1$ and in $W_\R$ by theta decay. The topology equivalence and
$t_\Gamma(h)\le C\norm h_{H^1}^2$ make them Cauchy in the original
energy topology. Subtracting a constant is an admissible gauge change,
preserving energy and variance.

For an exterior limit with zero moment, choose even compact smooth
$\psi$ in $O$ with $m_O(\psi)=1$. If $u_n$ are killed-core approximants,
replace them by $u_n-m_O(u_n)\psi$. Moment continuity makes the
correction vanish in form norm. Hence the centered cut lies in
precisely $\D_O^0$, not just in the moment kernel of a larger domain.

